% Additional h10 experiment, based on the full supplied record.
\subsection{Additional Evaluation on Undulating Terrain}
\label{subsec:undulating_h10}

The OFAT-Optimized configuration obtained in the preceding stages was
further evaluated on an undulating circular track (Fig.~\ref{fig:h10_setup}).
The terrain height was 10\,mm from trough to crest (h10), with a sinusoidal
amplitude of 5\,mm and a wavelength of 0.60\,m. A single native heightfield
was simulated at an 8\,ms time step. Terrain-contact settings were friction
1, restitution 0, soft CFM $10^{-5}$, and a maximum of 10 contact joints;
these differ from the default-contact settings used for the flat-track trials.

\begin{figure}[htbp]
  \centering
  \includegraphics[width=\linewidth]{Figure/H10_Setup.png}
  \caption{VoBiRo on the undulating circular track used for the additional
  h10 evaluation. The black line provides the visual path reference.}
  \label{fig:h10_setup}
\end{figure}

\begin{figure}[!htbp]
  \centering
  \begin{subfigure}[b]{0.48\linewidth}
    \centering
    \includegraphics[width=\linewidth]{Figure/H10_Roll_Panel.pdf}
    \subcaption{}
    \label{fig:h10_roll}
  \end{subfigure}
  \hfill
  \begin{subfigure}[b]{0.48\linewidth}
    \centering
    \includegraphics[width=\linewidth]{Figure/H10_Pitch_Panel.pdf}
    \subcaption{}
    \label{fig:h10_pitch}
  \end{subfigure}
  \caption{Kinematic responses on the h10 undulating track over the first
  11\,s: (a) CoM Roll representing lateral sway and (b) CoM Pitch
  representing sagittal tilt.}
  \label{fig:h10_responses}
\end{figure}

Figure~\ref{fig:h10_responses} presents CoM Roll and CoM Pitch over the
first 11\,s, using the same display window as the preceding comparisons.
These terms follow the earlier sections and denote IMU-derived orientation
angles. Full-record metrics were computed from all 13538 samples over
108.296\,s, giving a mean GPS speed of 0.1756\,m/s and zero-reference
CoM Roll and CoM Pitch RMSE values of 0.1699 and 0.1345\,rad, respectively.
The mean pitch was $-0.1118$\,rad, so its RMSE includes a sustained posture
offset as well as oscillation.

No physical fall was observed during the recording. The logger's
instantaneous acceleration-threshold flags were not treated as verified
falls, and no samples were excluded. This trial extends the evaluation
to the tested undulating-terrain and contact configuration. With only
one h10 recording and no matched Default trial on the same terrain,
the results do not establish comparative improvement or trial-to-trial
robustness.

\FloatBarrier

